% from ega3-1-fr.tex
\oldpage[III]{85}
Ces derniers vérifient évidemment la condition de transitivité
$\varphi_{\mathbf f^{m-p}}=
\varphi_{\mathbf f^{m-n}}\circ\varphi_{\mathbf f^{n-p}}$ pour $p\leq n\leq m$~;
ils définissent donc deux \emph{systèmes inductifs} de $A$-modules~; on posera
\begin{align}
\label{III.1.1.6.4-fr}
\tag{1.1.6.4}
C^\bullet((\mathbf f),M)&=\varinjlim_n K^\bullet(\mathbf f^n,M),\\
\label{III.1.1.6.5-fr}
\tag{1.1.6.5}
H^\bullet((\mathbf f),M)&=H^\bullet(C^\bullet((\mathbf f),M))
=\varinjlim_n H^\bullet(\mathbf f^n,M),
\end{align}
la dernière égalité provenant du fait que le passage à la limite inductive
permute au foncteur $H^\bullet$ (G, I, 2.1). On verra ultérieurement
(\hyperref[III.1.4.3-fr]{1.4.3}) que $H^\bullet((\mathbf f),M)$ ne dépend en
fait que de l'\emph{idéal $(\mathbf f)$} dans $A$ (et même de la topologie
$(\mathbf f)$-préadique sur $A$), ce qui justifie les notations.

Il est clair que $M\mapsto C^\bullet((\mathbf f),M)$ est un foncteur
$A$-linéaire exact, et $M\mapsto H^\bullet((\mathbf f),M)$ un foncteur
cohomologique.
\end{env}

\begin{env}[1.1.7]
\label{III.1.1.7-fr}
Soient $\mathbf f=(f_i)\in A^r$, $\mathbf g=(g_i)\in A^r$~; désignons par
$e_{\mathbf g}$ la multiplication à gauche par le vecteur
$\mathbf g\in A^r$ dans l'algèbre extérieure $\bigwedge(A^r)$~; on sait que
l'on a la \emph{formule d'homotopie}
\begin{equation}
\label{III.1.1.7.1-fr}
\tag{1.1.7.1}
i_{\mathbf f}e_{\mathbf g}+e_{\mathbf g}i_{\mathbf f}
=\langle\mathbf g,\mathbf f\rangle 1
\end{equation}
dans le $A$-module $A^r$ ($1$ désignant l'automorphisme identique de $A^r$)~;
cette relation signifie aussi que, \emph{dans le complexe
$K_\bullet(\mathbf f)$}, on a
\begin{equation}
\label{III.1.1.7.2-fr}
\tag{1.1.7.2}
de_{\mathbf g}+e_{\mathbf g}d=\langle\mathbf g,\mathbf f\rangle 1.
\end{equation}
Si l'idéal $(\mathbf f)$ est égal à $A$, il existe $\mathbf g\in A^r$ tel que
$\langle\mathbf g,\mathbf f\rangle=\sum_{i=1}^r g_if_i=1$. Par suite
(G, I, 2.4)~:
\end{env}

\begin{proposition}[1.1.8]
\label{III.1.1.8-fr}
Supposons que l'idéal $(\mathbf f)$ engendré par les $f_i$ soit égal à $A$.
Alors le complexe $K_\bullet(\mathbf f)$ est homotopiquement trivial, et il en
est donc de même des complexes $K_\bullet(\mathbf f,M)$ et
$K^\bullet(\mathbf f,M)$ pour tout $A$-module $M$.
\end{proposition}

\begin{corollary}[1.1.9]
\label{III.1.1.9-fr}
Si $(\mathbf f)=A$, on a $H^\bullet(\mathbf f,M)=0$ et
$H^\bullet((\mathbf f),M)=0$ pour tout $A$-module $M$.
\end{corollary}

\begin{proof}
En effet, on a alors $(\mathbf f^n)=A$ pour tout $n$.
\end{proof}

\begin{remark}[1.1.10]
\label{III.1.1.10-fr}
Avec les mêmes notations que ci-dessus, soient $X=\operatorname{Spec}(A)$,
$Y$ le sous-schéma fermé de $X$ défini par l'idéal $(\mathbf f)$. Nous
prouverons au \S~9 que $H^\bullet((\mathbf f),M)$ est isomorphe à la
cohomologie $H_Y^\bullet(X,\widetilde M)$ correspondant à l'antifiltre $\Phi$
des parties fermées de $Y$ (T, 3.2). Nous montrerons aussi que la prop.
(\hyperref[III.1.2.3-fr]{1.2.3}) appliquée à $X$ et à
$\mathcal F=\widetilde M$, est un cas particulier d'une suite exacte de
cohomologie
\[
\cdots\longrightarrow H_Y^p(X,\mathcal F)\longrightarrow H^p(X,\mathcal F)
\longrightarrow H^p(X-Y,\mathcal F)\longrightarrow H_Y^{p+1}(X,\mathcal F)
\longrightarrow\cdots.
\]
\end{remark}

\subsection{Cohomologie de Čech d'un recouvrement ouvert}
\label{subsection:III.1.2-fr}

\begin{notation}[1.2.1]
\label{III.1.2.1-fr}
Dans ce numéro, on notera~:
\begin{itemize}
\item $X$ un préschéma~;
\item $\mathcal F$ un $\mathcal O_X$-Module quasi-cohérent~;
